\documentclass[11pt]{article}
\usepackage[utf8]{inputenc}
\usepackage[T1]{fontenc}
\usepackage{amsmath,amssymb,mathtools}
\usepackage[margin=1in]{geometry}
\usepackage{hyperref}
\usepackage{array}
\usepackage{parskip}
\setlength{\parindent}{0pt}
\hypersetup{colorlinks=true,linkcolor=blue,urlcolor=blue}
\title{Lecture 10: Convergence, Properties, and DTFS}
\author{}
\date{}
\begin{document}
\maketitle

\section*{Dirichlet Conditions for CTFS Convergence}

The Fourier series converges for signals satisfying all three:

\begin{enumerate}
  \item \textbf{Absolute integrability:} $\int_T |x(t)| \, dt < \infty$
  \item \textbf{Bounded variation:} Finite number of maxima/minima per period
  \item \textbf{Finite discontinuities:} Finite number of finite-jump discontinuities per period
\end{enumerate}

At a \textbf{discontinuity}, the FS converges to the midpoint:

\[x_{FS}(t_0) = \frac{x(t_0^-) + x(t_0^+)}{2}\]

\section*{Gibbs Phenomenon}

\begin{itemize}
  \item Near jump discontinuities, finite partial sums exhibit \textbf{overshoot ripples}
  \item Peak overshoot is always \textasciitilde{}\textbf{9\%} of jump height, regardless of number of harmonics
  \item Overshoot region narrows as $N \to \infty$ but peak amplitude stays at \textasciitilde{}9\%
  \item A finite sum of smooth sinusoids \textbf{cannot} reproduce a discontinuity
\end{itemize}

\section*{Smoothness and Spectral Decay}

\begin{center}
\begin{tabular}{|l|l|l|}
\hline
\textbf{Signal Type} & \textbf{Spectral Decay} & \textbf{Smoothness} \\
\hline
Jump discontinuity & $\|a_k\| \sim 1/k$ & 0 continuous derivatives \\
\hline
Continuous, $dx/dt$ discontinuous & $\|a_k\| \sim 1/k^2$ & 1st derivative discontinuous \\
\hline
General: $m$th derivative discontinuous & $\|a_k\| \sim 1/k^{m+1}$ & $m$ continuous derivatives \\
\hline
\end{tabular}
\end{center}

\textbf{Key insight:} Smoother signals need fewer harmonics for accurate representation.

\section*{CTFS Properties}

\begin{center}
\begin{tabular}{|l|l|l|}
\hline
\textbf{Property} & \textbf{Time Domain} & \textbf{Frequency Domain} \\
\hline
Linearity & $Ax(t) + By(t)$ & $Aa_k + Bb_k$ \\
\hline
Time Shift & $x(t - t_0)$ & $a_k \, e^{-jk\omega_0 t_0}$ \\
\hline
Time Reversal & $x(-t)$ & $a_{-k}$ \\
\hline
Conjugation & $x^*(t)$ & $a_{-k}^*$ \\
\hline
Differentiation & $dx/dt$ & $jk\omega_0 \, a_k$ \\
\hline
Conjugate Symmetry & $x(t)$ real & $a_{-k} = a_k^*$ \\
\hline
\end{tabular}
\end{center}

\subsection*{Parseval's Relation (Average Power)}

\[P = \frac{1}{T} \int_T |x(t)|^2 \, dt = \sum_{k=-\infty}^{\infty} |a_k|^2\]

Power is conserved across domains. Each $|a_k|^2$ gives the power in the $k$th harmonic.

\section*{Key Property Implications}

\begin{itemize}
  \item \textbf{Time shift} changes phase but NOT magnitude: $|b_k| = |a_k|$, so power is preserved
  \item \textbf{Differentiation} amplifies high frequencies by factor $k$: acts as a highpass operation
  \item \textbf{DC component} ($k=0$) is always killed by differentiation: $d_0 = j(0)\omega_0 a_0 = 0$
\end{itemize}

\section*{Discrete-Time Fourier Series (DTFS)}

\subsection*{Key Differences from CTFS}

\begin{center}
\begin{tabular}{|l|l|l|}
\hline
\textbf{Feature} & \textbf{CTFS} & \textbf{DTFS} \\
\hline
Period & $T$ (continuous) & $N$ (integer samples) \\
\hline
Fundamental freq & $\omega_0 = 2\pi/T$ & $\omega_0 = 2\pi/N$ \\
\hline
Sum range & $k = -\infty$ to $\infty$ & $k = 0$ to $N-1$ (finite!) \\
\hline
Convergence issues & Yes (Dirichlet/Gibbs) & \textbf{None} (finite sum) \\
\hline
Coefficients & Aperiodic & \textbf{Periodic}: $a_{k+N} = a_k$ \\
\hline
\end{tabular}
\end{center}

\subsection*{DTFS Pair}

\textbf{Synthesis:}
\[x[n] = \sum_{k=\langle N \rangle} a_k \, e^{jk(2\pi/N)n}\]

\textbf{Analysis:}
\[a_k = \frac{1}{N} \sum_{n=\langle N \rangle} x[n] \, e^{-jk(2\pi/N)n}\]

\subsection*{Why Only N Coefficients?}

Because $e^{j(k+N)(2\pi/N)n} = e^{jk(2\pi/N)n} \cdot e^{j2\pi n} = e^{jk(2\pi/N)n}$

The $(k+N)$th harmonic is \textbf{identical} to the $k$th harmonic. Adding more harmonics beyond $N$ produces no new signals.

\subsection*{DT Parseval's Relation}

\[\frac{1}{N} \sum_{n=\langle N \rangle} |x[n]|^2 = \sum_{k=\langle N \rangle} |a_k|^2\]

\end{document}
